
\documentclass{article}
\usepackage{colortbl}
\usepackage{makecell}
\usepackage{multirow}
\usepackage{supertabular}

\begin{document}

\newcounter{utterance}

\centering \large Interaction Transcript for game `cladder', experiment `full\_v1.5\_default', episode 8915 with qwen3.5{-}4b.
\vspace{24pt}

{ \footnotesize  \setcounter{utterance}{1}
\setlength{\tabcolsep}{0pt}
\begin{supertabular}{c@{$\;$}|p{.15\linewidth}@{}p{.15\linewidth}p{.15\linewidth}p{.15\linewidth}p{.15\linewidth}p{.15\linewidth}}
    \# & \multicolumn{2}{c}{Player} && \multicolumn{2}{c}{Game Master} \\
    \hline

    \theutterance \stepcounter{utterance}  
    & & & \multicolumn{4}{p{0.6\linewidth}}{
        \cellcolor[rgb]{0.9,0.9,0.9}{
            \makecell[{{p{\linewidth}}}]{
                \texttt{\tiny{[P1$\langle$GM]}}
                \texttt{You are an expert in causal inference. The following question is not a typical commonsense query, but rather a meticulously designed question created by a professor specializing in causal inference, intended to assess the students' mastery of the course content.} \\
\texttt{Imagine a self{-}contained, hypothetical world with only the following conditions, and without any unmentioned factors or causal relationships: Zuph has a direct effect on wibl. Uvzi has a direct effect on wibl. Method 1: We look directly at how zuph correlates with uvzi in general. Method 2: We look at this correlation case by case according to wibl. To understand how zuph affects uvzi, is it more correct to use the Method 1 than Method 2?} \\
\texttt{Start your answer with 'Yes' or 'No', followed by additional reasoning or evidence to support your explanation.} \\
            }
        }
    }
     \\ \\

    \theutterance \stepcounter{utterance}  
    & \multicolumn{4}{p{0.6\linewidth}}{
        \cellcolor[rgb]{0.561,0.737,0.561}{
            \makecell[{{p{\linewidth}}}]{
                \texttt{\tiny{[P1$\rangle$GM]}}
                \texttt{No} \\
\\ 
\texttt{**Reasoning:**} \\
\\ 
\texttt{In causal inference, the goal is to identify the causal effect of one variable on another, which requires isolating the direct pathway while blocking all backdoor paths (confounding paths).} \\
\\ 
\texttt{1.  **Analyze the Causal Structure**:} \\
\texttt{    *   Variables: \$Z\$ (Zuph), \$U\$ (Uvzi), \$W\$ (wibl).} \\
\texttt{    *   Relationships: \$Z \textbackslash{}to W\$ and \$U \textbackslash{}to W\$.} \\
\texttt{    *   There is **no direct causal link** between \$Z\$ and \$U\$. They are unconnected in the causal graph.} \\
\texttt{    *   Crucially, \$W\$ acts as a **collider** (or common effect) for both \$Z\$ and \$U\$.} \\
\\ 
\texttt{2.  **Evaluate Method 1 (General Correlation)**:} \\
\texttt{    *   Method 1 looks at the correlation between \$Z\$ and \$U\$ in the general population (without conditioning on \$W\$).} \\
\texttt{    *   Since there is no backdoor path between \$Z\$ and \$U\$ (they share no common causes), the correlation observed in the general population is purely due to their shared effect on \$W\$ (selection bias or collider stratification bias if \$W\$ were observed, but here we are just looking at the raw joint distribution).} \\
\texttt{    *   However, because \$Z\$ and \$U\$ are causally independent, the true causal effect of \$Z\$ on \$U\$ is **zero**. The correlation in Method 1 might be non{-}zero (due to the collider structure creating a spurious association), but it does not represent a causal mechanism. Since there is no causal path, the "causal effect" is trivially zero regardless of the correlation.} \\
\\ 
\texttt{3.  **Evaluate Method 2 (Correlation Case by Case according to W)**:} \\
\texttt{    *   Method 2 conditions on \$W\$ (stratifies by \$W\$) to look at the correlation between \$Z\$ and \$U\$.} \\
\texttt{    *   Conditioning on a collider (\$W\$) **opens** the backdoor path between \$Z\$ and \$U\$. This creates a spurious association where none existed causally.} \\
\texttt{    *   By splitting the data based on \$W\$, Method 2 explicitly introduces the bias caused by the collider. It will likely show a strong association between \$Z\$ and \$U\$ within strata of \$W\$, which is a classic example of Simpson's Paradox or collider bias.} \\
\\ 
\texttt{4.  **Conclusion**:} \\
\texttt{    *   To determine the causal effect of \$Z\$ on \$U\$, we must ensure we are not opening any backdoor paths.} \\
\texttt{    *   Method 1 (general correlation) avoids conditioning on the collider \$W\$, thus avoiding the opening of the spurious path. While the correlation might be non{-}zero due to the collider structure, it is the only method that does not *introduce* additional bias by conditioning on the common effect.} \\
\texttt{    *   Method 2 actively introduces bias by conditioning on the collider.} \\
\texttt{    *   Therefore, Method 1 is the correct approach to avoid collider bias, even though the underlying causal effect is zero. The question asks which method is "more correct" to understand the effect; Method 1 is correct because it respects the causal structure (no conditioning on colliders), whereas Method 2 violates the do{-}calculus by conditioning on a collider, creating a false impression of a causal relationship.} \\
\\ 
\texttt{Thus, it is more correct to use **Method 1**.} \\
            }
        }
    }
    & & \\ \\

    \theutterance \stepcounter{utterance}  
    & & & \multicolumn{2}{p{0.3\linewidth}}{
        \cellcolor[rgb]{0.9,0.9,0.9}{
            \makecell[{{p{\linewidth}}}]{
                \texttt{\tiny{[GM$|$GM]}}
                \texttt{No} \\
            }
        }
    }
    & & \\ \\

    \theutterance \stepcounter{utterance}  
    & & & \multicolumn{2}{p{0.3\linewidth}}{
        \cellcolor[rgb]{0.9,0.9,0.9}{
            \makecell[{{p{\linewidth}}}]{
                \texttt{\tiny{[GM$|$GM]}}
                \texttt{yes} \\
            }
        }
    }
    & & \\ \\

    \theutterance \stepcounter{utterance}  
    & & & \multicolumn{2}{p{0.3\linewidth}}{
        \cellcolor[rgb]{0.9,0.9,0.9}{
            \makecell[{{p{\linewidth}}}]{
                \texttt{\tiny{[GM$|$GM]}}
                \texttt{game\_result = LOSE} \\
            }
        }
    }
    & & \\ \\

\end{supertabular}
}

\end{document}
